\documentclass[10.5pt]{article}
\usepackage{fontspec}\defaultfontfeatures{Ligatures=TeX}
\usepackage[letterpaper,margin=0.6in]{geometry}
\usepackage{amsmath,xcolor,longtable,array,tikz}
\usetikzlibrary{arrows.meta}
\setlength{\parindent}{0pt}\setlength{\parskip}{4pt}
\definecolor{pblue}{HTML}{1971C2}\definecolor{pred}{HTML}{E03131}\definecolor{pgreen}{HTML}{2F9E44}\definecolor{ppurple}{HTML}{7048E8}\definecolor{pgray}{HTML}{868E96}
\newfontfamily\mono{TeX Gyre Cursor}
\newcommand{\mtag}[2]{\colorbox{#1}{\textcolor{white}{\textsf{\fontsize{6.5}{7}\selectfont\bfseries #2}}}}
\newcommand{\clock}[1]{{\mono\fontsize{7}{8}\selectfont #1}\ }
% numbered step: \step{num}{tagcolor}{TAG}{text}
\newcommand{\step}[4]{\par\noindent\makebox[0.32in][l]{\textbf{\small #1}}\mtag{#2}{#3}\ {\small #4}\par\vspace{2pt}}
\newcommand{\pic}[1]{\par\hspace{0.32in}\begin{tikzpicture}[scale=0.42,every node/.style={font=\scriptsize}]#1\end{tikzpicture}\par\vspace{3pt}}
\newcommand{\CHECK}[3]{\par\noindent\makebox[0.32in][l]{\textbf{\small #1}}\fcolorbox{black}{black!4}{\parbox{2.55in}{\textsf{\fontsize{6.5}{8}\selectfont\bfseries CHECK}\ {\fontsize{7.5}{9}\selectfont #2}\par{\fontsize{7.5}{9}\selectfont\textbf{Answer:} #3}\par{\fontsize{7}{8.5}\selectfont\color{pgray} 4 of 5 right: go on. Fewer: one more example, ask again.}}}\par\vspace{3pt}}
\newcommand{\WATCH}[1]{\par\hspace{0.32in}{\textsf{\fontsize{6.5}{8}\selectfont\bfseries\color{pred} WATCH FOR}\ {\small\itshape\color{pred} #1}}\par}
\newcommand{\bl}[1]{\underline{\textcolor{pgreen}{\textbf{#1}}}}
\newcommand{\pen}[2]{\textcolor{#1}{\rule{8pt}{8pt}}\ #2}
\newcommand{\picL}[1]{\par\hspace{0.32in}\begin{tikzpicture}[scale=0.6,every node/.style={font=\footnotesize}]#1\end{tikzpicture}\par\vspace{3pt}}
% shell picture: x, y, shells (list of electron counts is drawn by \shells), label
\newcommand{\ring}[4]{\draw[pblue] (#1,#2) circle (#3);\foreach \k in {1,...,#4}{\pgfmathsetmacro{\ang}{90+360*\k/#4}\fill ({#1+#3*cos(\ang)},{#2+#3*sin(\ang)}) circle (0.07);}}
\newcommand{\nuc}[3]{\fill[pred!20] (#1,#2) circle (0.3);\node[pred] at (#1,#2){\tiny #3};}
% Lewis dots: \sd{x}{y}{n}{offset x}{offset y}{pair spread x}{pair spread y}
\newcommand{\sd}[7]{\ifnum#3=1 \fill ({#1+#4},{#2+#5}) circle (0.07);\fi\ifnum#3=2 \fill ({#1+#4-#6},{#2+#5-#7}) circle (0.07);\fill ({#1+#4+#6},{#2+#5+#7}) circle (0.07);\fi}
% \lw{x}{y}{symbol}{top}{right}{bottom}{left}: a box with the symbol and the dots
\newcommand{\lw}[7]{\draw[pgray] (#1-0.65,#2-0.65) rectangle (#1+0.65,#2+0.65);\node at (#1,#2){\scriptsize #3};\sd{#1}{#2}{#4}{0}{0.45}{0.14}{0}\sd{#1}{#2}{#5}{0.45}{0}{0}{0.14}\sd{#1}{#2}{#6}{0}{-0.45}{0.14}{0}\sd{#1}{#2}{#7}{-0.45}{0}{0}{0.14}}
% a bare box with a number in it
\newcommand{\nb}[3]{\draw[pgray] (#1-0.65,#2-0.65) rectangle (#1+0.65,#2+0.65);\node at (#1,#2){\scriptsize #3};}
\pagestyle{plain}
\begin{document}
\begin{center}{\Large\bfseries Unit 4, Day 2 board script: Why Atoms Bond and Lewis Dot Symbols}\\[2pt]{\normalsize College prep \textbullet{} 50 minutes \textbullet{} \texttt{chemistry/whiteboard/index.html\#lesson=bonding}}\end{center}
\noindent\textbf{Pens.} \pen{pblue}{shells, distance} \quad \pen{pred}{protons, pull} \quad \pen{pgreen}{new words (labels)} \quad \pen{ppurple}{shared electrons} \quad \pen{black}{everything else, including every valence dot}\par
\noindent\textbf{Steps.} Numbered by frame: 1.1, 1.2\ldots{} are on frame 1. Each DRAW step has a small picture of what the board should look like when the step is done. Tick the green Words box as you write each label.\par
\noindent\textbf{Dot rule.} One dot on each side (top, right, bottom, left) first; then pair up, starting at the top. This is the same rule the unit test uses.\par
\noindent\textbf{Not today.} No ions or charges (day 4), no shared pairs (day 7), no energy language, no electron configurations, no predicting how many bonds an atom makes. The reason atoms bond is always ``to get a full outer level, 8 valence electrons.''\par\smallskip
\noindent\fcolorbox{black}{black!3}{\parbox{\dimexpr\linewidth-2\fboxsep}{\small
\textbf{Time checkpoints.} Be starting frame 2 by \textbf{14:00}, frame 3 by \textbf{22:00}, frame 4 by \textbf{30:00}, the exit check by \textbf{38:00}. The checks are never cut; trim around them.\par
\textbf{If you are behind,} cut in this order: (1) step 4.4, the worked nitrogen: ask it as a quick whiteboard question instead; (2) step 1.7, ``period = shells'': say it in one sentence, draw nothing; (3) step 4.6, the preview: one sentence; (4) exit question b; (5) the summary becomes tonight's homework. If you reach 38:00 still in frame 4, finish check 4, skip to the exit check, and give the summary as homework.\par
\textbf{If you are ahead} (the 4-of-5 checks pass quickly), use these before the exit check, one at a time on whiteboards: \emph{Draw the Lewis dot symbol for sulfur.} (Top pair, right pair, one dot each on bottom and left.) \emph{How many unpaired electrons does selenium have?} (2: group 16, 8 $-$ 6.) \emph{Nitrogen or fluorine: which has more unpaired electrons?} (Nitrogen, 3 to 1.) \emph{Argon or chlorine: which one bonds? Why?} (Chlorine: argon's outer level is already full.)}}\par\medskip
\begin{longtable}{@{}>{\setlength{\parskip}{5pt}}p{3.55in}@{\hspace{0.25in}}p{3.05in}@{}}
\multicolumn{2}{@{}l}{\large\bfseries Frame 0: question and warm-up \hfill{\mono\small 0:00--7:00}}\tabularnewline*[2pt]
\textbf{Essential question.} Why do atoms bond, and what do the dots in a Lewis symbol show?\par
\textbf{Warm-up.} Shells in sulfur? Sodium or potassium: more reactive? Why are noble gases the exception? &
\step{0.1}{black}{SAY}{\clock{0:00}Warm-up is up as they walk in. At 3:00 remove the covers: 3; potassium; their outer shell is already full.}
\step{0.2}{black}{SAY}{\clock{5:00}Read the question. Point at the color key: ``Blue is always shells. Red is always protons. Today the dots are black.''} \tabularnewline[10pt]
\multicolumn{2}{@{}l}{\large\bfseries Frame 1: valence electrons \hfill{\mono\small 7:00--14:00}}\tabularnewline*[2pt]
The electrons in an atom's outermost shell are its \bl{valence} electrons. Atoms in the same \bl{group} have the same number of them.\par
For groups 13 to 18, subtract \bl{10} from the group number. Helium is the exception: it has only \bl{2}. The period number is the number of \bl{shells}, not the number of valence electrons.\par\smallskip
\textbf{1.} A student says sulfur has 3 valence electrons because it is in period 3. What is wrong with that? &
\step{1.1}{pblue}{DRAW}{In the empty space: sodium, nucleus $+11$ and three rings. Put 2, 8 and 1 electrons on the rings.}
\pic{\ring{0}{0}{0.55}{2}\ring{0}{0}{1.1}{8}\ring{0}{0}{1.65}{1}\nuc{0}{0}{+11}\node at (0,-2.1){Na};}
\step{1.2}{pgreen}{DRAW}{Circle the outer electron. Write VALENCE beside it. Say: ``The outermost shell does the bonding.''}
\pic{\ring{0}{0}{0.55}{2}\ring{0}{0}{1.1}{8}\ring{0}{0}{1.65}{1}\nuc{0}{0}{+11}\node at (0,-2.1){Na};\draw[pgreen,thick] (0,1.65) circle (0.28);\draw[pgreen,-{Stealth}] (2.6,1.9)--(0.4,1.7);\node[pgreen,right] at (2.6,1.9){VALENCE};}
\step{1.3}{black!60}{ASK}{``Sodium is in group 1. How many valence electrons do lithium and potassium have?'' (One each: same group.)}
\tabularnewline
&
\step{1.4}{pgreen}{DRAW}{In the row of boxes under the table: 1, 2, then 3 to 8, one box per group. Write GROUP above the first column.}
\pic{\foreach \i/\n in {0/1,1/2,2/3,3/4,4/5,5/6,6/7,7/8}{\nb{\i*1.4}{0}{\n}}\node[pgreen] at (0,1.1){GROUP};}
\step{1.5}{pgreen}{WORDS}{Tick \textbf{valence} and \textbf{group}.}
\step{1.6}{pgreen}{DRAW}{Beside the row: ``13 to 18: group $-$ 10'' and ``He: 2.'' Tick \textbf{10} and \textbf{2}.}
\pic{\node[pgreen,right] at (0,0){13 to 18: group $-$ \textbf{10}};\node[pgreen,right] at (0,-1){He: \textbf{2}};}
\step{1.7}{pblue}{DRAW}{Down the left of the table: ``PERIOD = SHELLS, not valence electrons.'' Tick \textbf{shells}.}
\pic{\node[pblue,right] at (0,0){PERIOD = \textcolor{pgreen}{SHELLS}};\node[right] at (0,-1){not valence electrons};}
\CHECK{1.8}{Sulfur has 3 valence electrons because it is in period 3?}{No: period 3 means 3 shells. Sulfur is in group 16, so 6 valence electrons.}
\WATCH{``Period 3, so 3 valence electrons.'' Ask: does the period count shells, or outer electrons?} \tabularnewline[10pt]
\multicolumn{2}{@{}l}{\large\bfseries Frame 2: the octet rule \hfill{\mono\small 14:00--22:00}}\tabularnewline*[2pt]
Noble gases already have a \bl{full} outer level: 8 valence electrons (helium has 2). That makes them very \bl{stable}, so they hardly bond.\par
Every other atom bonds to get a full outer level. This is the \bl{octet} rule.\par
\textbf{Example.} Sodium has 1 valence electron, so it \bl{loses} it and is left with a full level, like neon. Chlorine has 7, so it \bl{gains} 1 to make 8.\par\smallskip
\textbf{2.} Neon has 8 valence electrons and fluorine has 7. Why does fluorine bond but neon does not? &
\step{2.1}{black!60}{ASK}{Point at the neon picture. ``How many valence electrons? Is the outer level full?'' (8; yes.)}
\step{2.2}{pgreen}{DRAW}{Under Ne: FULL, STABLE, no bond.}
\pic{\ring{0}{0}{0.55}{2}\ring{0}{0}{1.1}{8}\nuc{0}{0}{+10}\node at (0,-1.6){Ne};\node[pgreen] at (0,-2.3){FULL};\node[pgreen] at (0,-3){STABLE};\node at (0,-3.7){no bond};}
\step{2.3}{black!60}{ASK}{``Sodium has 1 and chlorine has 7. How far is each from a full outer level?'' Let two answers come before you draw.}
\tabularnewline
&
\step{2.4}{pblue}{DRAW}{Under Na: an arrow carrying its outer electron away. Write LOSES 1.}
\pic{\ring{0}{0}{0.55}{2}\ring{0}{0}{1.1}{8}\ring{0}{0}{1.65}{1}\nuc{0}{0}{+11}\node at (0,-2.1){Na};\draw[pblue,-{Stealth}] (0.2,1.8)--(2.3,1.8);\node[pgreen] at (0,-2.9){LOSES 1};}
\step{2.5}{pred}{DRAW}{Under Cl: an arrow bringing one electron in to the outer ring. Write GAINS 1.}
\pic{\ring{0}{0}{0.55}{2}\ring{0}{0}{1.1}{8}\ring{0}{0}{1.65}{7}\nuc{0}{0}{+17}\node at (0,-2.1){Cl};\draw[pred,-{Stealth}] (2.4,1.3)--(1.5,0.9);\node[pgreen] at (0,-2.9){GAINS 1};}
\step{2.6}{pgreen}{DRAW}{Between the pictures: OCTET RULE, and ``8 valence electrons = full outer level.''}
\pic{\node[pgreen] at (0,0){OCTET RULE};\node at (0,-0.8){8 valence electrons = full outer level};}
\step{2.7}{pgreen}{WORDS}{Tick \textbf{full}, \textbf{stable}, \textbf{octet}, \textbf{loses}, \textbf{gains}.}
\CHECK{2.8}{Neon has 8 valence electrons and fluorine has 7. Why does fluorine bond but neon does not?}{Neon's outer level is already full; fluorine is one electron short of an octet.}
\WATCH{``Atoms bond because they want to.'' Say it the same way each time: to get a full outer level.} \tabularnewline[10pt]
\multicolumn{2}{@{}l}{\large\bfseries Frame 3: Lewis dot symbols \hfill{\mono\small 22:00--30:00}}\tabularnewline*[2pt]
A Lewis dot symbol shows only the \bl{valence} electrons. The symbol stands for the nucleus and the inner electrons.\par
Put one dot on each of the four sides (top, right, bottom, left) before you \bl{pair} any dots. Then pair up, starting at the top.\par
\textbf{Example.} Nitrogen is in group 15, so it has \bl{5} valence electrons: one dot on each side, and the fifth pairs with the \bl{top} dot.\par\smallskip
\textbf{3.} A student draws carbon with two pairs of dots, one on top and one on the right. Which rule did they break? &
\step{3.1}{black}{SAY}{``The letters are the nucleus and the inner electrons. Only the valence electrons become dots.'' Tick \textbf{valence}.}
\step{3.2}{black}{DRAW}{In the boxes: Li, one dot on top. Be, second dot on the right. B, third on the bottom. C, fourth on the left. Say: ``One on each side first.''}
\picL{\lw{0}{0}{Li}{1}{0}{0}{0}\lw{1.9}{0}{Be}{1}{1}{0}{0}\lw{3.8}{0}{B}{1}{1}{1}{0}\lw{5.7}{0}{C}{1}{1}{1}{1}}
\step{3.3}{pgreen}{DRAW}{N, O, F, Ne: the fifth dot pairs with the top one, the sixth with the right one, and so on. Write PAIR under the row.}
\picL{\lw{0}{0}{N}{2}{1}{1}{1}\lw{1.9}{0}{O}{2}{2}{1}{1}\lw{3.8}{0}{F}{2}{2}{2}{1}\lw{5.7}{0}{Ne}{2}{2}{2}{2}\node[pgreen] at (2.85,-1.2){PAIR};}
\tabularnewline
&
\step{3.4}{black!60}{ASK}{``Count the dots in each box. What do you notice about the group?'' (The dots match the valence electrons.)}
\step{3.5}{black}{DRAW}{In the empty space, nitrogen from scratch. Say each dot: ``Five valence electrons. Four singles, one on each side. The fifth pairs with the top.'' Tick \textbf{5} and \textbf{top}.}
\picL{\lw{0}{0}{N}{2}{1}{1}{1}\node[pgreen,right] at (1,0.4){5 valence};\node[pgreen,right] at (1,-0.4){fifth pairs on top};}
\CHECK{3.6}{Carbon drawn with two pairs, one on top and one on the right: which rule was broken?}{One dot on each side before any pair forms. Carbon has four single dots.}
\WATCH{Counting the inner electrons (``sodium has 11 dots'') or pairing before every side has a dot.} \tabularnewline[10pt]
\multicolumn{2}{@{}l}{\large\bfseries Frame 4: unpaired electrons \hfill{\mono\small 30:00--38:00}}\tabularnewline*[2pt]
An \bl{unpaired} electron is a dot that is alone on its side. Count the single dots.\par
Up to group 14, every valence electron is unpaired. From group 15 on, unpaired electrons $=$ \bl{8} minus the valence electrons. Noble gases have \bl{0} unpaired electrons, so there is nothing to pair up and they do not bond.\par
\textbf{Example.} Nitrogen has 5 valence electrons, so $8-5=$ 3 unpaired electrons.\par\smallskip
\textbf{4.} Oxygen is in group 16. How many unpaired electrons does it have? Show how you know. &
\step{4.1}{black}{SAY}{``A dot alone on its side is an unpaired electron.'' Tick \textbf{unpaired}.}
\step{4.2}{pgreen}{DRAW}{In the boxes under the groups, write how many single dots each atom has: 1, 2, 3, 4, 3, 2, 1, 0. Point back at frame 3 and count with them.}
\pic{\foreach \i/\n in {0/1,1/2,2/3,3/4,4/3,5/2,6/1,7/0}{\nb{\i*1.4}{0}{\n}}\node[pgreen,left] at (-0.8,0){UNPAIRED};}
\step{4.3}{black!60}{ASK}{``What happens to the count after group 14?'' (It goes back down.) ``Why?'' (The extra dots are pairing up.)}
\tabularnewline
&
\step{4.4}{black}{DRAW}{Next to ``Worked: nitrogen'': 5 valence, $8-5=3$ unpaired. Think aloud: find the group, count the valence electrons, take it from 8. Tick \textbf{8}.}
\pic{\node[right] at (0,0){N: 5 valence};\node[right] at (0,-0.9){8 $-$ 5 = \textcolor{pgreen}{3 UNPAIRED}};}
\step{4.5}{pgreen}{DRAW}{Next to ``Noble gases'': 0 UNPAIRED, nothing to pair up. Ask: ``How does that match the octet rule?'' (Full outer level, no reason to bond.) Tick \textbf{0}.}
\pic{\node[right] at (0,0){Noble gases: \textcolor{pgreen}{0}};\node[right] at (0,-0.9){nothing to pair up};}
\step{4.6}{black}{SAY}{One sentence for the next lessons: ``Atoms use their unpaired electrons to bond: they give them away, take more, or share them.'' Draw nothing.}
\CHECK{4.7}{Oxygen is in group 16. How many unpaired electrons? Show how you know.}{6 valence electrons, so $8-6=2$ unpaired.}
\WATCH{``Oxygen has 6 unpaired electrons.'' Ask: how many dots are alone?} \tabularnewline[10pt]
\multicolumn{2}{@{}l}{\large\bfseries Frame 5: exit check and summary \hfill{\mono\small 38:00--50:00}}\tabularnewline*[2pt]
\textbf{a.} Draw the Lewis dot symbol for magnesium and for chlorine. \textbf{b.} How many unpaired electrons does silicon have? \textbf{c.} Why does argon not form bonds? One sentence.\par
\textbf{Summary.} Answer the essential question in two or three sentences. Use the words ``octet'' and ``valence.'' &
\step{5.1}{black}{SAY}{\clock{38:00}Alone, silent, 6 minutes. Walk the room. Remove the covers at 44:00: Mg two single dots (top, right); Cl pairs top, right and bottom with one dot on the left; 4; argon has a full outer level of 8 and no unpaired electrons.}
\step{5.2}{black}{SAY}{\clock{45:00}Summary: silent writing, read two aloud. Listen for ``octet'' for why atoms bond and ``valence'' for what the dots show.}
\step{5.3}{black}{SAY}{\clock{48:00}``Sodium has one valence electron and chlorine has seven. What do you think happens when they meet?'' Stop there.} \tabularnewline
\end{longtable}
\end{document}
